\begin{abstract}
Information-driven follow-up planners for astronomical transients assume that the event belongs to a modeled class and that measurements can be taken whenever they are most informative. We test a planner that drops both assumptions. It keeps an explicit unmodeled hypothesis $h_0$, represented by a Gaussian process, in its posterior, and adds an opportunity-decay term to its utility. Three preregistered hypotheses were tested in a simulated constellation of eight Starlink satellites with real orbital visibility, with measurements strictly forward in time (\SimNEvents\ events per scenario). At a 5\% false-alarm rate, the planner's $P(h_0)$ detected \SimTprP\ of events from withheld classes, against \SimTprBthree\ for a classifier-entropy threshold, while closed-world information gain was confidently wrong on \SimCwWhBthreeAll\ of them. Removing the decay term reduced recovered amplitude information more for fast classes than for slow ones, although the effect was carried by kilonovae, reversed for flares and reversed under model error. The planner was not the fastest to identify modeled events: classifier-entropy sampling was. An adequacy term in the utility contributed nothing measurable. On \RealNEvents\ real ZTF transients replayed against the same fleet, a preregistered causal test confirmed four of five hypotheses: the planner was confidently wrong on \RealCCwWhP\ of real tidal disruption events, superluminous and SN~Ibn supernovae, against \RealCCwWhBthree\ for closed-world information gain, with an AUROC of \RealCAuc\ and supernova accuracy equal to the closed-world planner's; it could not be shown non-inferior in accuracy within 5 points. We also report a time-ordering error that affected our earlier results, and how correcting it changed them.
\end{abstract}
